\scanpage{216}
Démonstration du lemme d. \deleted{Une section de $X$ sur $S$ est définie par un hom \deleted{\unsure{d'algèbres}} de $k[t]$-algèbres graduées $\varphi:B\to A[z]$ ; \unsure{si on peut} \unsure{on peut} prendre $\varphi(z)=z$, $\varphi(x)=\xi$, \deleted{\unsure{$\varphi(y)=\eta$}} avec $\xi,\eta\in A$, $\xi,\eta$ sont soumis à $\xi\eta-t$ D'ailleurs} On a un ouvert \inserted{affine} $V$ de $X$ isomorphe à $\operatorname{Spec} k[x,y,t]/(xy-t)$, et une section de $V$ sur $S$ est un \inserted{$k[t]$-}hom $k[x,y,t]/(xy-t)\longrightarrow k[t]$, i.e. donné par deux éléments $\xi,\eta$ de $k[t]$ tels que $\xi\eta-t=0$

1) $\xi=1$, $\eta=t$ définit la section $s_1$. On regarde pour la spécialisation $t=0$

2) $\xi=t$, $\eta=1$ \enspace ---

\deleted{Proposition} Démonstration du corollaire. \deleted{Soit \uline{Lemme e} Sous les conditions, de $D$ un diviseur sur $X$. \unsure{Sa} restriction à $X\mid U$ est $\simeq$ à un diviseur $D'_1$ [car} On a \deleted{\unsure{et} $H^1(U,\unsure{A|L}_U^*)$}
\[
\begin{aligned}
H^1(X\mid U,\underline{\mathcal O}_{X\mid U}^*)&=H^1(P^1_U,\underline{\mathcal O}_{P^1_U}^*)\simeq\\
H^1(U,\underline{\mathcal O}_U^*)\times\deleted{\unsure{H^1}}\mathbb Z&\simeq\mathbb Z
\end{aligned}
\]
compte tenu que $H^1(U,\underline{\mathcal O}_U^*)=0$ puisque $k[t][t^{-1}]$ est principal

[ N.B. on n'a pas besoin du th. \unsure{entier|entièrement} pénible sur $H^1(S,P^1_S)$, \deleted{\unsure{S étant}} qui ici se simplifie beaucoup grâce au fait que $S$ est noeth\unsure{ér.|.}]
% archive sheet number: 215
